
\section{Pipeline}

%\note{following LDM-151}
\todo{consider a (simplified?) pipeline graph}

\subsection{Template Generation}

%\note{Pre-DR1 this may require description independent of DRP.}
%
%High-level discussion of DCR correction could go here.

\subsection{Preload}

\subsection{Single Frame Processing}

%\note{Construction of calibration products assumed described in \cite{PSTN-026}.}

ISR, PSF and Background fitting, photometric and astrometric calibration.  Likely major commonalities with DRP.

\subsection{Image Differencing}

\subsection{Source Detection and Measurement}

Including discussion of point source, dipole, and streak detection on difference images.  

%\note{Algorithmic basics could be discussed elsewhere}

\subsection{Initial Filtering} \label{sec:filtering}

\texttt{FilterDiaSourceCatalogTask} in \texttt{lsst.ap.association}

\subsection{Reliability scoring}

Image-differencing searches for transients typically contend with high rates of false positives.
Among the raw detections, artifacts may dominate real astrophysical sources by an order of magnitude.
These may be due to imperfectly-corrected instrument signatures, astrometric mismatches between the template and the science image, cosmic rays, or algorithmic failures in the differencing.
Modern surveys employ machine-learned classifiers to winnow these candidates \citep[e.g.,][]{Bailey:07:NearbySNFML,Bloom:12:RealBogus,Brink:13:RB2,Wright:15:PS1ML,Goldstein:15:DESRealBogus,Duev:19:Braii}.
These have steadily improved in performance and sophistication as the classifiers transitioned from Random Forest models trained on manually-constructed feature sets to deep neural networks trained directly on image pixels.
Some approaches have combined both detection and scoring without using the difference image \citep{Sedaghat:17:Transinet,Acero-Cuellar:22:NoDifference}.

\todo{overview of the algorithm, training, and performance of the ML spuriousness score; detailed discussion likely deferred to separate paper}

\subsection{Catalog Transformation}

\texttt{TransformDiaSourceCatalogTask} in  \texttt{lsst.ap.association}

\eric{schema tables?}

\subsection{Source Association}

Due to AP's real-time nature, it is not possible to perform \textit{post-facto} source association as in the annual data releases.
Instead, spatial association is performed on-the-fly by the \texttt{DiaPipelineTask} within the \texttt{lsst.ap.association} package.
A dedicated Alert Production Database (APDB; \S \ref{sec:db}) holds the current state of the system.

During the preload step (\S \ref{sec:preload}), DIAObjects and SSObjects overlapping the expected field of view of the image are stored in the local prompt processing worker butler repository.
The new DIASources are first spatially associated with the existing DIAObjects by finding the closest match within a maximum distance of one arcsecond.
Pairs with the closest spatial separations are joined first to minimize misassociations.
When a match is found, the DIASource is added to the DIAObject's list of measurements.
Unassociated DIASources are then spatially associated with the predicted positions of the SSObjects at the time of the exposure.
For DIASources with no matching DIAObject or SSObject, a new DIAObject is created and the DIASource is added to it.



%\note{does \cite{PSTN-045} describe the APDB/PPDB implementation?}

\subsection{Alert Generation}

Each DIASource not filtered for quality purposes (\S \ref{sec:filtering}) produces an alert.
To enable downstream users to quickly identify events of interest, the alert packets contain a range of information.  
They include:

\begin{itemize}
    \item The triggering DIASource record
    \item The corresponding DIAObject or SSObject record
    \item Any DIASource and DIAForcedSource records that exist, and difference image noise estimates where they do not, taken from the previous 12 months.
    \item Cutout images of the science, template, and difference images.
\end{itemize}
\eric{confirm contents}

The cutout images are packaged in the FITS format using the astropy \texttt{CCDData} class and include mask and variance planes, an estimate of the PSF, a World Coordinate System, and other metadata.
Most cutouts are 30$\times$30 pixels but are expanded to a larger size for extended sources.
The maximimum cutout size of \todo{confirm} allows cutouts to contain trailed Near Earth Asteroids traveling slower than ten degrees per day.

Alerts are transmitted and stored in Apache Avro\footnote{\url{https://avro.apache.org/}} format, a strongly-schemaed compact binary serialization.
However, the Avro standard includes its schemas as uncompressed ASCII.
To conserve bandwidth, we therefore send alerts without including the schemas required to deserialize them.
Instead, the schemas are stored in a separate schema registry and are referenced by a unique identifier in the alert packet.

\eric{sizing estimates here or in initial performance section?}

\subsection{Alert Distribution} \label{sec:alertdist}

We use the Kafka distributed streaming platform\footnote{\url{https://kafka.apache.org/}} to transmit alerts to brokers.
Kafka provides a fault-tolerant, low-overhead mechanism to transmit large volumes of data to multiple consumers and is widely used both in industry and in astronomy. 
\eric{cite patterson, GCN, Hopskotch, other Rubin usage...}

The Alert Distribution Kafka system is deployed at the USDF using Kubernetes.
\bri{details of the deployment}
\eric{schema management}

Because of the large size of the alert packets and the short latency requirement, bandwidth out of the US Data Facility (USDF) at SLAC limits the number of direct recepients of the full alert stream.
Seven community alert brokers were selected through a proposal process \eric{cite} to receive the full alert stream directly from the USDF:
\eric{list}
Two additional brokers will operate downstream of a full-stream broker.


\subsection{Alert Filtering Service}

ANTARES system

community alert filters

\subsection{Forced Photometry}

\subsection{Source Injection}

\subsection{Metrics and Monitoring}